\textbf{结论名称：} 高次不等式解法（穿根法/穿针引线法，奇穿偶不穿）

\textbf{适用条件：} 一元高次不等式能分解为一次因式幂的乘积，且不等式为 $f(x)>0$、$f(x)<0$ 等严格或非严格形式。通常先将最高次项系数化为正数。

\textbf{变量说明：} 设 $f(x)=a(x-x_1)^{m_1}(x-x_2)^{m_2}\cdots(x-x_n)^{m_n}$，其中 $a>0$，各零点从小到大排列 $x_1<x_2<\cdots<x_n$，$m_i\in\mathbb{N}^*$ 为因式的重数（幂次）。

\textbf{核心公式：} 
\[
\boxed{\text{奇穿偶不穿：在数轴上标记各根，从右上方起画曲线，遇到 } m_i \text{ 为奇数的根穿过 } x \text{ 轴，} m_i \text{ 为偶数的根不穿过}}
\]

\textbf{使用场景：} 快速求解形如 $(x-1)(x+2)^2(x-3)^3>0$ 的一元高次不等式，避免分段讨论，直接根据曲线在 $x$ 轴上、下的位置写出解集。

\textbf{简单例子：} 解不等式 $(x-1)(x+2)^2(x-3)^3>0$。\\
根为 $-2$（偶，$m=2$），$1$（奇，$m=1$），$3$（奇，$m=3$）。\\
从右上角起：过 $x=3$ 穿至轴下，过 $x=1$ 穿回轴上，过 $x=-2$ 不穿始终在轴上。\\
曲线在 $x$ 轴上方的区间为 $(-\infty,-2)$、$(-2,1)$、$(3,+\infty)$，且 $x=-2$ 使因式为零应剔除，故解集为 $(-\infty,-2)\cup(-2,1)\cup(3,+\infty)$。